% Resonance duck mode parameter sweep results table. Companion to
% parameter_sweep_table.tex (Advanced) / parameter_sweep_table_basic.tex
% (Basic) - same color-coding convention, applied to Resonance's 100-point
% sweep instead of the original 9-run manual sweep.
%
% Data source: analysis/run_resonance_sweep.py --n 100 (Construction Scene,
% Dialogue priority). Resonance mode reuses the same Sidechain Compressor
% law (threshold/ratio/knee/attack/release) as Basic/Advanced, applied per
% spectral peak, plus two Resonance-only parameters: num_peaks (how many
% spectral peaks are tracked) and bandwidth_octaves (how wide each peak's
% gain-reduction bump is) - see ResonanceSuppressor.h/Engine.h.
%
% METHODOLOGY NOTE: this sweep required two real fixes before these numbers
% were trustworthy. (1) main.cpp's --static export mode did not support
% --mode resonance at all before this sweep was built - added --peaks/
% --bandwidth/--resonance-max-reduction flags. (2) Resonance mode has a
% fixed 2048-sample (42.7ms @ 48kHz) algorithmic latency that neither
% compute_stoi()'s lag search (+/-20ms) nor compute_masking_margin() (no
% lag correction at all) accounted for - the first run of this sweep
% produced nonsensical numbers (margin gain exceeding anything Basic/
% Advanced ever produced, STOI/loudness uniformly and severely negative)
% purely from comparing temporally-misaligned signals. Confirmed via a
% wide +/-100ms cross-correlation search independently landing on exactly
% 2048 samples; fixed by trimming that known, deterministic latency from
% the rendered mix before scoring. See analysis/resonance_sweep_findings.txt
% for the full writeup this table summarizes.
%
% Requires packages already in main.tex's preamble: booktabs, siunitx (the
% footnote below uses \SI{}{}). Also requires colortbl for \cellcolor - see
% parameter_sweep_table.tex's header comment (add \usepackage{colortbl} or
% change to \usepackage[table]{xcolor}).
%
% Plain right-aligned (r) columns throughout, not siunitx S columns - same
% reasoning as parameter_sweep_table.tex: \cellcolor combined with an
% S column's number-parsing is fragile, so precision is fixed by hand
% instead.
%
% COLOR CODING: the 7 swept-input columns (Thr./Ratio/Knee/Atk./Rel./
% Peaks/Bandwidth) get grey header backgrounds; Margin Gain/STOI Gain/
% Loud. Red. are shaded per-cell, intensity scaled linearly between this
% table's own min and max (i.e. relative to these 5 curated rows only, not
% the full 100-point sweep or any other table's scale).
%
% ROW SELECTION: not a ranked top-N (unlike parameter_sweep_table.tex) -
% these 5 rows are chosen specifically to show the sweep's two headline
% findings side by side: (1) the best point found regardless of peak count,
% plus the best 1-peak and best 2-peak points, are all close to each other
% and all use wide bandwidth - peak count doesn't separate them; (2) the
% "Narrow bandwidth" row uses a comparably aggressive threshold and a
% moderate peak count (4) but a narrow bandwidth (0.12 octaves) and
% collapses to almost no effect - isolating bandwidth, not threshold or
% peak count, as the reason it fails; (3) the "Weak threshold" row is the
% same near-zero floor every other sweep in this project shows once
% threshold is too conservative to engage at all.

\definecolor{paramgrey}{gray}{0.85}

\begin{table}[h]
\centering
\tiny
\caption{Resonance duck mode parameter sweep (100-point Latin Hypercube,
Construction Scene, Dialogue priority): five representative points
isolating bandwidth, not peak count, as the parameter that separates a
strong result from a weak one. Gains are relative to the unprocessed
baseline. Best point found here, ranked by STOI gain (margin +3.67dB,
STOI +0.070, loudness reduction +10.6pp), falls meaningfully short of
Advanced mode's own 150-point-sweep ceiling on the same scene (margin
+6.96dB, STOI +0.104) -- loudness reduction is the one metric where
Resonance is competitive (this sweep's own loudness-reduction ceiling,
+11.5pp, essentially matches Advanced's +11.4pp).}
\label{tab:resonance-sweep}
\begin{tabular}{@{}l r r r r r r r r r r@{}}
\toprule
\cellcolor{paramgrey}\textbf{Point} & \cellcolor{paramgrey}\textbf{Thr.} & \cellcolor{paramgrey}\textbf{Ratio} & \cellcolor{paramgrey}\textbf{Knee} & \cellcolor{paramgrey}\textbf{Atk.} & \cellcolor{paramgrey}\textbf{Rel.} & \cellcolor{paramgrey}\textbf{Peaks} & \cellcolor{paramgrey}\textbf{BW (oct)} & \textbf{Margin Gain} & \textbf{STOI Gain} & \textbf{Loud.\ Red.} \\
 & \cellcolor{paramgrey}(dB) & \cellcolor{paramgrey}{} & \cellcolor{paramgrey}(dB) & \cellcolor{paramgrey}(ms) & \cellcolor{paramgrey}(ms) & \cellcolor{paramgrey}{} & \cellcolor{paramgrey}{} & (dB) & & (pp) \\
\midrule
Best overall (3 peaks)   & -50.6 & 2.4 & 11.4 & 46.0 & 319 & 3 & 3.98 & \cellcolor{green!70}+3.67 & \cellcolor{yellow!70}+0.070 & \cellcolor{orange!66}+10.6 \\
Best 1-peak              & -56.0 & 6.1 & 5.0  & 13.9 & 230 & 1 & 3.56 & \cellcolor{green!50}+2.36 & \cellcolor{yellow!66}+0.065 & \cellcolor{orange!70}+11.5 \\
Best 2-peak              & -52.4 & 8.8 & 23.6 & 16.8 & 381 & 2 & 2.28 & \cellcolor{green!41}+1.76 & \cellcolor{yellow!46}+0.040 & \cellcolor{orange!54}+8.1  \\
Narrow bandwidth (contrast) & -54.7 & 9.3 & 6.5  & 24.3 & 155 & 4 & 0.12 & \cellcolor{green!16}+0.08 & \cellcolor{yellow!20}+0.006 & \cellcolor{orange!20}+1.0  \\
Weak threshold (floor)   & -18.2 & 6.0 & 13.7 & 29.3 & 174 & 2 & 2.73 & \cellcolor{green!15}+0.00 & \cellcolor{yellow!15}+0.000 & \cellcolor{orange!15}+0.0  \\
\bottomrule
\end{tabular}
\par\vspace{4pt}
\raggedright\footnotesize Note the "Narrow bandwidth" row uses a threshold
about as aggressive as the top three rows (\SI{-54.7}{\decibel}) and a
moderate peak count (4) -- what collapses its result to almost nothing is
specifically its narrow \SI{0.12}{octave} bandwidth, not its threshold or
peak count. Across the full 100-point sweep, bandwidth is the
second-strongest predictor of STOI gain after threshold (standardized
regression coefficient $+0.505$, vs.\ threshold's $-0.775$, $R^2=0.813$),
while peak count's independent contribution is negligible ($+0.061$) --
1--2 peaks and 3--6 peaks perform statistically indistinguishably once
bandwidth is already wide. Margin gain is a weaker signal for Resonance
mode specifically than for Basic/Advanced: it is measured on Dialogue's
fixed nominal band, but Resonance ducks wherever the key's spectral peaks
currently are, not a fixed band -- see
\texttt{analysis/resonance\_sweep\_findings.txt} for the full caveat.
\end{table}
